
\documentclass{article}
\usepackage{colortbl}
\usepackage{makecell}
\usepackage{multirow}
\usepackage{supertabular}

\begin{document}

\newcounter{utterance}

\centering \large Interaction Transcript for game `dond', experiment `coop\_en', episode 14 with reflexion{-}Qwen3.5{-}27B.
\vspace{24pt}

{ \footnotesize  \setcounter{utterance}{1}
\setlength{\tabcolsep}{0pt}
\begin{supertabular}{c@{$\;$}|p{.15\linewidth}@{}p{.15\linewidth}p{.15\linewidth}p{.15\linewidth}p{.15\linewidth}p{.15\linewidth}}
    \# & $\;$A & \multicolumn{4}{c}{Game Master} & $\;\:$B\\
    \hline

    \theutterance \stepcounter{utterance}  
    & & \multicolumn{4}{p{0.6\linewidth}}{
        \cellcolor[rgb]{0.9,0.9,0.9}{
            \makecell[{{p{\linewidth}}}]{
                \texttt{\tiny{[P1$\langle$GM]}}
                \texttt{You are playing a cooperative negotiation game in which you have to agree on how to divide a set of items among you and another player.} \\
\\ 
\texttt{Rules:} \\
\texttt{(a) You and the other player are given a set of items. Each of you is also given a secret value function, representing how much you value each type of object.} \\
\texttt{(b) You exchange messages with the other player to agree on who gets which items. You can send a maximum of 5 messages each, or terminate early by making a secret proposal at any time.} \\
\texttt{(c) You are each asked to submit a secret proposal indicating the items you want formatted in square brackets as follows: "{[}Proposal: <number> <object name>, <number> <object name>, <...>{]}"} \\
\texttt{(d) If your proposals are complementary, i.e., there are enough items to fulfill both proposals, each player is awarded a score based on the sum of values for the items they received. Otherwise, both of you get zero points.} \\
\texttt{(e) Your goal is to maximize the sum of your score and the score of the other player.} \\
\\ 
\texttt{Let us start.} \\
\\ 
\texttt{The set of available items is:} \\
\\ 
\texttt{3 rings, 1 doll, 2 rocks} \\
\\ 
\texttt{Your secret value function is:} \\
\\ 
\texttt{ring: 2, doll: 0, rock: 2} \\
\\ 
\texttt{IMPORTANT: Your messages, unless it is the secret proposal, are directly transmitted to the other player, so do not include any response to the rules or text announcing your message. To make a secret proposal, use the indicated format. Do not use square brackets when communicating to the other player or it will be interpreted as your secret proposal.} \\
            }
        }
    }
    & \\ \\

    \theutterance \stepcounter{utterance}  
    & \multicolumn{4}{p{0.6\linewidth}}{
        \cellcolor[rgb]{0.9,0.9,0.9}{
            \makecell[{{p{\linewidth}}}]{
                \texttt{\tiny{[P1$\rangle$GM]}}
                \texttt{I value rings and rocks equally at 2 points each, and I have no use for the doll (0 points). Since there are 3 rings and 2 rocks available, and I need to maximize our combined score, I suggest we split the items based on our preferences.} \\
\\ 
\texttt{Could you share your value function for the items? That way, we can figure out who should take the doll and how to best divide the rings and rocks to ensure we both get high scores.} \\
            }
        }
    }
    & & \\ \\

    \theutterance \stepcounter{utterance}  
    & & & \multicolumn{2}{p{0.3\linewidth}}{
        \cellcolor[rgb]{0.9,0.9,0.9}{
            \makecell[{{p{\linewidth}}}]{
                \texttt{\tiny{[GM$|$GM]}}
                \texttt{continue} \\
            }
        }
    }
    & & \\ \\

    \theutterance \stepcounter{utterance}  
    & & & \multicolumn{2}{p{0.3\linewidth}}{
        \cellcolor[rgb]{0.9,0.9,0.9}{
            \makecell[{{p{\linewidth}}}]{
                \texttt{\tiny{[GM$|$GM]}}
                \texttt{I value rings and rocks equally at 2 points each, and I have no use for the doll (0 points). Since there are 3 rings and 2 rocks available, and I need to maximize our combined score, I suggest we split the items based on our preferences.} \\
\\ 
\texttt{Could you share your value function for the items? That way, we can figure out who should take the doll and how to best divide the rings and rocks to ensure we both get high scores.} \\
            }
        }
    }
    & & \\ \\

    \theutterance \stepcounter{utterance}  
    & & \multicolumn{4}{p{0.6\linewidth}}{
        \cellcolor[rgb]{0.9,0.9,0.9}{
            \makecell[{{p{\linewidth}}}]{
                \texttt{\tiny{[GM$\rangle$P2]}}
                \texttt{You are playing a cooperative negotiation game in which you have to agree on how to divide a set of items among you and another player.} \\
\\ 
\texttt{Rules:} \\
\texttt{(a) You and the other player are given a set of items. Each of you is also given a secret value function, representing how much you value each type of object.} \\
\texttt{(b) You exchange messages with the other player to agree on who gets which items. You can send a maximum of 5 messages each, or terminate early by making a secret proposal at any time.} \\
\texttt{(c) You are each asked to submit a secret proposal indicating the items you want formatted in square brackets as follows: "{[}Proposal: <number> <object name>, <number> <object name>, <...>{]}"} \\
\texttt{(d) If your proposals are complementary, i.e., there are enough items to fulfill both proposals, each player is awarded a score based on the sum of values for the items they received. Otherwise, both of you get zero points.} \\
\texttt{(e) Your goal is to maximize the sum of your score and the score of the other player.} \\
\\ 
\texttt{Let us start.} \\
\\ 
\texttt{The set of available items is:} \\
\\ 
\texttt{3 rings, 1 doll, 2 rocks} \\
\\ 
\texttt{Your secret value function is:} \\
\\ 
\texttt{ring: 0, doll: 6, rock: 2} \\
\\ 
\texttt{IMPORTANT: Your messages, unless it is the secret proposal, are directly transmitted to the other player, so do not include any response to the rules or text announcing your message. To make a secret proposal, use the indicated format. Do not use square brackets when communicating to the other player or it will be interpreted as your secret proposal.} \\
\\ 
\\ 
\texttt{I value rings and rocks equally at 2 points each, and I have no use for the doll (0 points). Since there are 3 rings and 2 rocks available, and I need to maximize our combined score, I suggest we split the items based on our preferences.} \\
\\ 
\texttt{Could you share your value function for the items? That way, we can figure out who should take the doll and how to best divide the rings and rocks to ensure we both get high scores.} \\
            }
        }
    }
    & \\ \\

    \theutterance \stepcounter{utterance}  
    & & & \multicolumn{4}{p{0.6\linewidth}}{
        \cellcolor[rgb]{0.9,0.9,0.9}{
            \makecell[{{p{\linewidth}}}]{
                \texttt{\tiny{[GM$\langle$P2]}}
                \texttt{I value the doll at 6 points, while rings are worth 0 to me and rocks are worth 2 each. Since you have no use for the doll and I value it highly, it makes sense for me to take the doll.} \\
\\ 
\texttt{For the remaining items, you value rings and rocks equally at 2 points each, and I value rocks at 2 points but rings at 0. To maximize our combined score, you should take all the rings (since I get 0 from them) and we can split the rocks or you can take them all depending on what works best for you, as long as I get the doll.} \\
\\ 
\texttt{How about this: I take the 1 doll, and you take the 3 rings and 2 rocks? That way, I get 6 points, and you get (3 * 2) + (2 * 2) = 10 points, for a total of 16. Does that work for you?} \\
            }
        }
    }
     \\ \\

    \theutterance \stepcounter{utterance}  
    & & & \multicolumn{2}{p{0.3\linewidth}}{
        \cellcolor[rgb]{0.9,0.9,0.9}{
            \makecell[{{p{\linewidth}}}]{
                \texttt{\tiny{[GM$|$GM]}}
                \texttt{continue} \\
            }
        }
    }
    & & \\ \\

    \theutterance \stepcounter{utterance}  
    & & & \multicolumn{2}{p{0.3\linewidth}}{
        \cellcolor[rgb]{0.9,0.9,0.9}{
            \makecell[{{p{\linewidth}}}]{
                \texttt{\tiny{[GM$|$GM]}}
                \texttt{I value the doll at 6 points, while rings are worth 0 to me and rocks are worth 2 each. Since you have no use for the doll and I value it highly, it makes sense for me to take the doll.} \\
\\ 
\texttt{For the remaining items, you value rings and rocks equally at 2 points each, and I value rocks at 2 points but rings at 0. To maximize our combined score, you should take all the rings (since I get 0 from them) and we can split the rocks or you can take them all depending on what works best for you, as long as I get the doll.} \\
\\ 
\texttt{How about this: I take the 1 doll, and you take the 3 rings and 2 rocks? That way, I get 6 points, and you get (3 * 2) + (2 * 2) = 10 points, for a total of 16. Does that work for you?} \\
            }
        }
    }
    & & \\ \\

    \theutterance \stepcounter{utterance}  
    & & \multicolumn{4}{p{0.6\linewidth}}{
        \cellcolor[rgb]{0.9,0.9,0.9}{
            \makecell[{{p{\linewidth}}}]{
                \texttt{\tiny{[P1$\langle$GM]}}
                \texttt{I value the doll at 6 points, while rings are worth 0 to me and rocks are worth 2 each. Since you have no use for the doll and I value it highly, it makes sense for me to take the doll.} \\
\\ 
\texttt{For the remaining items, you value rings and rocks equally at 2 points each, and I value rocks at 2 points but rings at 0. To maximize our combined score, you should take all the rings (since I get 0 from them) and we can split the rocks or you can take them all depending on what works best for you, as long as I get the doll.} \\
\\ 
\texttt{How about this: I take the 1 doll, and you take the 3 rings and 2 rocks? That way, I get 6 points, and you get (3 * 2) + (2 * 2) = 10 points, for a total of 16. Does that work for you?} \\
            }
        }
    }
    & \\ \\

    \theutterance \stepcounter{utterance}  
    & \multicolumn{4}{p{0.6\linewidth}}{
        \cellcolor[rgb]{0.9,0.9,0.9}{
            \makecell[{{p{\linewidth}}}]{
                \texttt{\tiny{[P1$\rangle$GM]}}
                \texttt{That sounds like a perfect plan. Since I get 0 value from the doll and you get 6, and I get 2 from every ring and rock while you get 0 from rings, this allocation maximizes our total score to 16.} \\
\\ 
\texttt{I agree to:} \\
\texttt{{-} You take the 1 doll.} \\
\texttt{{-} I take the 3 rings and 2 rocks.} \\
\\ 
\texttt{{[}Proposal: 3 rings, 2 rocks{]}} \\
            }
        }
    }
    & & \\ \\

    \theutterance \stepcounter{utterance}  
    & & & \multicolumn{2}{p{0.3\linewidth}}{
        \cellcolor[rgb]{0.9,0.9,0.9}{
            \makecell[{{p{\linewidth}}}]{
                \texttt{\tiny{[GM$|$GM]}}
                \texttt{proposal} \\
            }
        }
    }
    & & \\ \\

    \theutterance \stepcounter{utterance}  
    & & & \multicolumn{2}{p{0.3\linewidth}}{
        \cellcolor[rgb]{0.9,0.9,0.9}{
            \makecell[{{p{\linewidth}}}]{
                \texttt{\tiny{[GM$|$GM]}}
                [3, 0, 2]
            }
        }
    }
    & & \\ \\

    \theutterance \stepcounter{utterance}  
    & & \multicolumn{4}{p{0.6\linewidth}}{
        \cellcolor[rgb]{0.9,0.9,0.9}{
            \makecell[{{p{\linewidth}}}]{
                \texttt{\tiny{[GM$\rangle$P2]}}
                \texttt{The other player has submitted a secret proposal. It is now your turn to submit a proposal of your own.} \\
            }
        }
    }
    & \\ \\

    \theutterance \stepcounter{utterance}  
    & & & \multicolumn{4}{p{0.6\linewidth}}{
        \cellcolor[rgb]{0.9,0.9,0.9}{
            \makecell[{{p{\linewidth}}}]{
                \texttt{\tiny{[GM$\langle$P2]}}
                \texttt{{[}Proposal: 3 rings, 2 rocks{]}} \\
            }
        }
    }
     \\ \\

    \theutterance \stepcounter{utterance}  
    & & & \multicolumn{2}{p{0.3\linewidth}}{
        \cellcolor[rgb]{0.9,0.9,0.9}{
            \makecell[{{p{\linewidth}}}]{
                \texttt{\tiny{[GM$|$GM]}}
                \texttt{proposal} \\
            }
        }
    }
    & & \\ \\

    \theutterance \stepcounter{utterance}  
    & & & \multicolumn{2}{p{0.3\linewidth}}{
        \cellcolor[rgb]{0.9,0.9,0.9}{
            \makecell[{{p{\linewidth}}}]{
                \texttt{\tiny{[GM$|$GM]}}
                [3, 0, 2]
            }
        }
    }
    & & \\ \\

    \theutterance \stepcounter{utterance}  
    & & & \multicolumn{2}{p{0.3\linewidth}}{
        \cellcolor[rgb]{0.9,0.9,0.9}{
            \makecell[{{p{\linewidth}}}]{
                \texttt{\tiny{[GM$|$GM]}}
                [[3, 0, 2], [3, 0, 2]]
            }
        }
    }
    & & \\ \\

    \theutterance \stepcounter{utterance}  
    & \multicolumn{4}{p{0.6\linewidth}}{
        \cellcolor[rgb]{0.95,0.95,0.95}{
            \makecell[{{p{\linewidth}}}]{
                \texttt{\tiny{[P2$\rightarrow$P1]}}
                \texttt{Post{-}episode Reflexion (not sent to GM):} \\
\texttt{New plan: I failed because I submitted a proposal that only listed the items I wanted, omitting the items the other player wanted. The rules state that proposals must be "complementary," meaning the sum of items in both proposals must equal the total available inventory. By only proposing "{[}3 rings, 2 rocks{]}", I left the doll unassigned in my proposal, causing a mismatch with the other player's likely proposal (which would include the doll) or simply failing to account for the full set. Next time, after agreeing on the split, I will ensure my secret proposal explicitly lists *only* the items I am claiming, but I must verify that the other player's proposal (which I cannot see) will cover the rest. However, the critical error is likely that the system expects the proposal to be a complete partition or that I misunderstood the "complementary" check. Actually, re{-}reading the rules: "If your proposals are complementary... there are enough items to fulfill both proposals." If I propose {[}3 rings, 2 rocks{]} and the other proposes {[}1 doll{]}, the union is the full set, so that should work. Wait, the error might be that I didn't wait for the other player to confirm their specific proposal format or I assumed the split was final without a formal counter{-}proposal. But the most likely strategic failure in these specific negotiation tasks is that I need to explicitly state the *entire* division in my proposal if the system interprets "complementary" as "my proposal + their proposal = total". If I only list my items, and they list theirs, it works. Why did it fail? Perhaps the other player didn't submit a proposal yet, or the system requires me to list the *agreed* split in a way that implies the other's share? No, the rule says "items you want".} \\
\\ 
\texttt{Let's re{-}evaluate the failure. The outcome was LOSE. The text says "every message you sent was an acceptable move". This implies the format was correct, but the logic failed. The logic was: I take rings/rocks, they take doll. I proposed {[}3 rings, 2 rocks{]}. If they proposed {[}1 doll{]}, it should be a win. Why did it lose?} \\
\texttt{Possibility A: The other player did not submit a proposal, and the game ended with my proposal being the only one, which is invalid.} \\
\texttt{Possibility B: The other player's proposal was different (e.g., they wanted a rock too), and I didn't negotiate that specific detail.} \\
\texttt{Possibility C: I misinterpreted the "complementary" rule. Maybe I need to list the *other* player's items in my proposal? No, "items you want".} \\
\texttt{Possibility D: The most common pitfall in these specific prompts is that the user (me) must ensure the *other* player actually submits a proposal. I sent a proposal immediately after agreeing.} \\
            }
        }
    }
    & & \\ \\

    \theutterance \stepcounter{utterance}  
    & & & \multicolumn{4}{p{0.6\linewidth}}{
        \cellcolor[rgb]{0.95,0.95,0.95}{
            \makecell[{{p{\linewidth}}}]{
                \texttt{\tiny{[P1$\rightarrow$P2]}}
                \texttt{Post{-}episode Reflexion (not sent to GM):} \\
\texttt{New plan: I failed because I proposed taking items (3 rings, 2 rocks) that I value at 0 and 2 respectively, while ignoring the doll which I value at 6. Although the other player likely took the doll, my proposal did not include it, meaning I received 0 points for the doll and the other player received 0 points for the rings/rocks they didn't get, or the proposals were not complementary in a way that maximized the sum. Specifically, I should have proposed taking the item I value most (the doll) and letting the other player take the items they value (rings and rocks). My new strategy is to explicitly include the doll in my proposal: "{[}Proposal: 1 doll{]}". This ensures I get my 6 points, and if the other player proposes taking the rings and rocks, our proposals will be complementary, maximizing the total score. I must align my proposal with my own value function, not just the other player's.} \\
            }
        }
    }
     \\ \\

\end{supertabular}
}

\end{document}
